\subsection{Zımmîlere Uygulanan Ceza Hukuku}
\indent\indent İslam dininde zorlama yoktur. Bu ilke doğrultusunda, dârülharbden ele geçirilen bölgeler sakinlerininin iki seçeneği bulunurdu. İsteyenler bölgeyi terk edip, dârülislam sınırları dışına çıkabilirlerdi. 
Kalanlar ise İslam hukukunun şartlarına uymayı ve askerlikten muaf kalmayı kabul ederek, cizye vergisi yükümlülüğünü üstlenirlerdi. Taraflar arasında yapıldığı varsayılan bu anlaşmaya zimmet anlaşması, kabul edenlere de zımmî denirdi.\footcite[293]{osmanli_hukuku}

Dârülislam sınırları içerisinde islam hukuku geçerliydi. Zımmîlerin davalarından da şer`î mahkemeler sorumluydu. Fakat Osmanlı Devleti, din ve vicdan hürriyeti kapsamında zımmîlere hukuki ve kazai muhtariyet 
tanımıştır. Bu sayede zımmîler kendi aralarındaki davaları cemaat mahkemelerinde görebiliyorlardı. Cemaat mahkemelerinin yaptırım gücü bulunmadığı için daha çok bir hakem kurulu işlevi görüyordu. Taraflar 
cemaat mahkemesinin verdiği hükmü kabul ederlerse, şer`î mahkemeler müdahil olmuyordu. Öte yandan taraflardan biri cemaat mahkemesinin kararını kabul etmezse, davayı şer`î mahkemelere taşıyabiliyordu. Şer`î 
mahkemeler, bu yapıda için itiraz mahkemesi işlevi görüyordu.

Cemaat mahkemelerine verilen imtiyazın ilk ne zaman ortaya çıktığı bilinmemektedir. Ancak bu imtiyazın sınırları zamanla daraltılmıştır. İslam hukukçuları zamanla bu imtiyazın sadece ahvâl-i şahsiyye dedikleri 
konularla sınırlı kalması gerektiği fikrini belirlemişlerdir. Bu yaklaşımın sonucu olarak, cemaat mahkemeleri sadece evlilik, boşanma, miras ve vasiyet gibi konularda yetkili kılınmıştır. Öte yandan da, cemaat 
mahkemelerinin zaman zaman ceza hukuku kapsamındaki davalara da baktığı bilinmektedir.\footcite[297]{osmanli_hukuku}

Şer`î mahkemelerde, gayrimüslimlere uygulanan ceza hukuku müslümanlara uygulanan ceza hukuku arasında bazı farklar bulunmaktaydı. Dinleri içki içmeye izin verdiği için kamu düzenini bozmadıkları sürece had cezası 
uygulanmamaktaydı. Zina suçunda ise, recm ya da had cezası yerine daha hafif bir ceza olan dayak cezası uygulanıyordu. Öte yandan domuz ve şarap gibi İslam hukukunda mütekavvim olmayan mallar, gayrimüslimler için 
mütekavvim sayılıyordu.\footcite[315]{osmanli_hukuku}

Gayrimüslimlere uygulanan bu imtiyazların, zaman içinde belirli grupların otonom bir yapı oluşturmalarına yol açtığı görülmektedir. Bu gruplara millet ismi verilmiştir. Osmanlı Devleti içindeki bütün ortodokslara
\textit{Rum milleti} denmiş ve Fener patrikhanesi bu milletin merkezi olmuştur. Benzer bir yapılanma \textit{Ermeni milleti} için de gerçekleşmiştir. Bu durum zaman içinde Osmanlı Devleti için sorun 
oluşturmaya başlamıştır. 1774 Küçük Kaynarca Anlaşması ile Osmanlı Devleti, dahilindeki ortodoksların din ve vicdan hürriyetinin korunması konusunda garanti vermiştir. Bu tarihten sonra gayrimülsimlerin durumları uluslarası bir hal almıştır.\footfullcite[148]{turk_hukuk_tarihi}